\documentclass[10pt]{article}
\usepackage[T1]{fontenc}
\usepackage[utf8]{inputenc}
\usepackage[french]{babel}
\usepackage{amsmath,amssymb,mathrsfs,mathpazo}
\usepackage{geometry}
\geometry{a4paper,margin=2cm}
\usepackage{graphicx,array,tabularx,tabularray,colortbl,multirow,multicol,booktabs,enumitem,pifont,ulem,fancybox,diagbox,cancel}
\usepackage{pstricks,pstricks-add,pst-plot,pst-tree,pst-node,pst-text,pst-func,pst-math,pst-bspline,pst-3dplot,pst-eucl,pst-solides3d}
\let\pscalculate\undefined
\usepackage{pst-bezier}
\usepackage{tikz,pgfplots,tkz-tab}
\usetikzlibrary{arrows,arrows.meta,patterns,calc,positioning,trees,decorations.markings,shapes.geometric,shapes.misc,shapes.arrows,intersections,decorations.pathmorphing,decorations.pathreplacing,backgrounds,shadows,babel}
\pgfplotsset{compat=1.18}
\usepgfplotslibrary{fillbetween}
\usepackage[np]{numprint}
\usepackage{listings,makecell,pas-tableur,twoopt}
\let\barre\undefined
\usepackage[pstricks]{tablvar}
\definecolor{darkgreen}{RGB}{0,100,0}
\providecommand{\red}{\color{red}}
\providecommand{\blue}{\color{blue}}
\providecommand{\green}{\color{green}}
\providecommand{\black}{\color{black}}
\providecommand{\white}{\color{white}}
\providecommand{\gray}{\color{gray}}
\providecommand{\cyan}{\color{cyan}}
\providecommand{\magenta}{\color{magenta}}
\providecommand{\yellow}{\color{yellow}}
\providecommand{\violet}{\color{violet}}
\providecommand{\orange}{\color{orange}}
\providecommand{\euro}{\texteuro}
\providecommand{\EUR}{\texteuro}
\providecommand{\ieme}{\textsuperscript{e}}
\providecommand{\degres}{\ensuremath{^\circ}}
\providecommand{\degree}{\ensuremath{^\circ}}
\providecommand{\vv}[1]{\overrightarrow{#1}}
\providecommand{\Coord}[1]{#1}
\providecommand{\decosix}{}
\providecommand{\starredbullet}{\ensuremath{\bullet}}
\pagestyle{empty}
\setlength{\parindent}{0pt}
\usepackage{adjustbox}\DeclareUnicodeCharacter{2236}{:}\DeclareUnicodeCharacter{2190}{\ensuremath{\leftarrow}}
\providecommand{\R}{\mathbb{R}}
\providecommand{\N}{\mathbb{N}}
\providecommand{\Z}{\mathbb{Z}}
\providecommand{\Q}{\mathbb{Q}}
\providecommand{\D}{\mathbb{D}}
\providecommand{\C}{\mathbb{C}}
\providecommand{\e}{{\mathrm{e}}}
\providecommand{\vect}[1]{\overrightarrow{#1}}
\providecommand{\vectt}[1]{\overrightarrow{\mathrm{#1}}}
\providecommand{\vvt}[1]{\overrightarrow{\mathrm{#1}}}
\providecommand{\np}[2][]{{#2~#1}}
\providecommand{\vv}[1]{{\overrightarrow{#1}}}
\renewcommand{\d}{{\mathrm{d}}}
\providecommand{\degres}{{^\circ}}
\providecommand{\degre}{{^\circ}}
\providecommand{\blue}{\color{blue}}
\providecommand{\red}{\color{red}}
\providecommand{\enskip}{\quad}
\providecommand{\llbracket}{\lbrack\!\lbrack}
\providecommand{\rrbracket}{\rbrack\!\rbrack}
\providecommand{\euro}{{\text{€}}}
\providecommand{\no}{{\text{n°}}}
\providecommand{\interval}[2]{{[#1;#2]}}
\providecommand{\intervFO}[2]{{[#1;#2[}}
\providecommand{\dotfill}{\ldots}

\providecommand{\ds}{\displaystyle}
\providecommand{\textding}[1]{#1}
\providecommand{\Oij}{\(\mathrm{O};\vec{\imath},\vec{\jmath}\)}
\providecommand{\Oijk}{\(\mathrm{O};\vec{\imath},\vec{\jmath},\vec{k}\)}
\providecommand{\Ouv}{\(\mathrm{O};\vec{u},\vec{v}\)}
\providecommand{\intervFF}[2]{\ensuremath{[#1;#2]}}
\renewcommand{\euro}{\ifmmode\text{\texteuro}\else\texteuro\fi}
\DeclareUnicodeCharacter{00B0}{\ifmmode^{\circ}\else\textdegree\fi}
\DeclareUnicodeCharacter{20AC}{\ifmmode\text{\texteuro}\else\texteuro\fi}
% Énoncés : Éducation nationale. Transcription et corrigés : APMEP, Denis Vergès et contributeurs. Découpage et classement : CoopMaths, Sébastien Lozano, Éric Elter et contributeurs. Adaptation de la présentation : Maths Explorer.
% CC BY-SA : https://coopmaths.fr/www/about/
% https://coopmaths.fr/alea/static/bac/2021/tex/bac_2021_09_metropole_5.tex
\begin{document}
\noindent\textit{APMEP / CoopMaths. CC BY-SA. Présentation adaptée par Maths Explorer.}\par\medskip
\medskip

\begin{tabular}{|l|}\hline
Principaux domaines abordés:\\
Fonction logarithme.\\ \hline
\end{tabular}

\bigskip

\textbf{Partie I}

\medskip

On considère la fonction $h$ définie sur l'intervalle $]0~;~ +\infty[$ par:

\[h(x) = 1 + \dfrac{\ln (x)}{x}.\]

\smallskip

\begin{enumerate}
\item Déterminer la limite de la fonction $h$ en $0$.
\item Déterminer la limite de la fonction $h$ en $+\infty$.
\item On note $h'$ la fonction dérivée de $h$. Démontrer que, pour tout nombre réel $x$ de $]0~;~ +\infty[$, on a:

\[h'(x) = \dfrac{1 - \ln (x)}{x^2}.\]

\smallskip

\item Dresser le tableau de variations de la fonction $h$ sur l'intervalle $]0~;~ +\infty[$.
\item Démontrer que l'équation $h(x) = 0$ admet une unique solution $\alpha$ dans $]0~;~ +\infty[$. 

Justifier que l'on a : $0,5 < \alpha < 0,6$.
\end{enumerate}

\bigskip

\textbf{Partie II}

\medskip

Dans cette partie, on considère les fonctions $f$ et $g$ définies sur $]0~;~ +\infty[$ par : 

\[f(x) = x \ln (x) - x ;\qquad  g(x) = \ln (x).\]

On note $\mathcal{C}_f$ et $\mathcal{C}_g$ les courbes représentant respectivement les fonctions $f$ et $g$ dans un repère orthonormé \Oij.

Pout tout nombre réel $a$ strictement positif, on appelle:

\setlength\parindent{1cm}
\begin{itemize}
\item[$\bullet~~$] $T_a$ la tangente à $\mathcal{C}_f$ en son point d'abscisse $a$ ;
\item[$\bullet~~$] $D_a$ la tangente à $\mathcal{C}_g$ en son point d'abscisse $a$.
\end{itemize}
\setlength\parindent{0cm}

Les courbes  $\mathcal{C}_f$ et $\mathcal{C}_g$ ainsi que deux tangentes $T_a$ et $D_a$ sont représentées ci-dessous.

\begin{center}
\psset{unit=1cm}
\begin{pspicture*}(-0.6,-2)(7,6)
\psgrid[gridlabels=0,subgriddiv=1,gridwidth=0.06pt]
\psaxes[linewidth=1.25pt,labelFontSize=\scriptstyle](0,0)(0,-2)(6.9,5.9)
\psaxes[linewidth=1.25pt,labelFontSize=\scriptstyle]{->}(0,0)(1,1)
\psplot[plotpoints=2000,linewidth=1.25pt,linecolor=red]{0.01}{7}{x ln x mul x sub}
\psplot[plotpoints=2000,linewidth=1.25pt,linecolor=blue]{0.01}{7}{x ln}
\psplot[plotpoints=2000,linewidth=1.25pt,linecolor=red,linestyle=dashed]{0.01}{7}{1.833 x mul 6.254 sub}
\psplot[plotpoints=2000,linewidth=1.25pt,linecolor=blue,linestyle=dashed]{0.01}{7}{0.162 x mul 0.833 add}
\uput[u](1,1){\blue $D_a$}\uput[r](2.9,-1){\red $T_a$}\uput[l](4.3,2){\red $\mathcal{C}_f$ }
\uput[r](0.2,-1.8){\blue $\mathcal{C}_g$}


\psline[linestyle=dotted,linewidth=1.25pt](6.25,0)(6.25,5.204)
\uput[d](6.25,0){$a$}
\end{pspicture*}
\end{center}

On recherche d'éventuelles valeurs de $a$ pour lesquelles les droites $T_a$ et $D_a$ sont perpendiculaires. 

Soit $a$ un nombre réel appartenant à l'intervalle $]0~;~ +\infty[$.


\medskip

\begin{enumerate}
\item Justifier que la droite $D_a$ a pour coefficient directeur $\dfrac{1}{a}$.
\item Justifier que la droite $T_a$ a pour coefficient directeur $\ln (a)$.
\end{enumerate}

On rappelle que dans un repère orthonormé, deux droites de coefficients directeurs respectifs $m$ et $m'$sont perpendiculaires si et seulement si $mm' = -1$.

\begin{enumerate}[resume]
\item Démontrer qu'il existe une unique valeur de $a$, que l'on identifiera, pour laquelle les droites $T_a$ et $D_a$ sont perpendiculaires.
\end{enumerate}

\end{document}
